%% Wrapper used by this directory's Makefile to export one TikZ figure as a
%% standalone, tight-cropped PDF. The book never compiles this -- it \inputs the
%% fragments directly via \tikzfig.
%%
%% The figure to draw is passed in on the command line:
%%     lualatex -jobname=NAME "\def\FIG{NAME.tex}%% Wrapper used by this directory's Makefile to export one TikZ figure as a
%% standalone, tight-cropped PDF. The book never compiles this -- it \inputs the
%% fragments directly via \tikzfig.
%%
%% The figure to draw is passed in on the command line:
%%     lualatex -jobname=NAME "\def\FIG{NAME.tex}%% Wrapper used by this directory's Makefile to export one TikZ figure as a
%% standalone, tight-cropped PDF. The book never compiles this -- it \inputs the
%% fragments directly via \tikzfig.
%%
%% The figure to draw is passed in on the command line:
%%     lualatex -jobname=NAME "\def\FIG{NAME.tex}%% Wrapper used by this directory's Makefile to export one TikZ figure as a
%% standalone, tight-cropped PDF. The book never compiles this -- it \inputs the
%% fragments directly via \tikzfig.
%%
%% The figure to draw is passed in on the command line:
%%     lualatex -jobname=NAME "\def\FIG{NAME.tex}\input{_wrapper}"
%%
%% Fonts match gwbook's defaults so an exported figure looks like what lands
%% on the page. Leading underscore keeps it out of the Makefile's wildcard.

\documentclass[border=6pt]{standalone}

\usepackage{fontspec}
\usepackage{unicode-math}
% unicode-math pulls amsmath in anyway; naming it makes the dependency explicit,
% since fragments use aligned/\dfrac and must typeset the same here as in the
% book, where gwbook.cls gets amsmath via physics.
\usepackage{amsmath}
\setmainfont{texgyreschola}[
  Extension=.otf, UprightFont=*-regular, ItalicFont=*-italic,
  BoldFont=*-bold, BoldItalicFont=*-bolditalic]
\setmathfont{texgyreschola-math.otf}

\usepackage{tikz}
% Same list, and for the same reason, as the preamble of gwbook.cls -- keep the
% two in step so a figure looks identical exported and on the page.
\usetikzlibrary{arrows.meta, calc, patterns.meta}

\begin{document}
\input{\FIG}
\end{document}
"
%%
%% Fonts match gwbook's defaults so an exported figure looks like what lands
%% on the page. Leading underscore keeps it out of the Makefile's wildcard.

\documentclass[border=6pt]{standalone}

\usepackage{fontspec}
\usepackage{unicode-math}
% unicode-math pulls amsmath in anyway; naming it makes the dependency explicit,
% since fragments use aligned/\dfrac and must typeset the same here as in the
% book, where gwbook.cls gets amsmath via physics.
\usepackage{amsmath}
\setmainfont{texgyreschola}[
  Extension=.otf, UprightFont=*-regular, ItalicFont=*-italic,
  BoldFont=*-bold, BoldItalicFont=*-bolditalic]
\setmathfont{texgyreschola-math.otf}

\usepackage{tikz}
% Same list, and for the same reason, as the preamble of gwbook.cls -- keep the
% two in step so a figure looks identical exported and on the page.
\usetikzlibrary{arrows.meta, calc, patterns.meta}

\begin{document}
\input{\FIG}
\end{document}
"
%%
%% Fonts match gwbook's defaults so an exported figure looks like what lands
%% on the page. Leading underscore keeps it out of the Makefile's wildcard.

\documentclass[border=6pt]{standalone}

\usepackage{fontspec}
\usepackage{unicode-math}
% unicode-math pulls amsmath in anyway; naming it makes the dependency explicit,
% since fragments use aligned/\dfrac and must typeset the same here as in the
% book, where gwbook.cls gets amsmath via physics.
\usepackage{amsmath}
\setmainfont{texgyreschola}[
  Extension=.otf, UprightFont=*-regular, ItalicFont=*-italic,
  BoldFont=*-bold, BoldItalicFont=*-bolditalic]
\setmathfont{texgyreschola-math.otf}

\usepackage{tikz}
% Same list, and for the same reason, as the preamble of gwbook.cls -- keep the
% two in step so a figure looks identical exported and on the page.
\usetikzlibrary{arrows.meta, calc, patterns.meta}

\begin{document}
\input{\FIG}
\end{document}
"
%%
%% Fonts match gwbook's defaults so an exported figure looks like what lands
%% on the page. Leading underscore keeps it out of the Makefile's wildcard.

\documentclass[border=6pt]{standalone}

\usepackage{fontspec}
\usepackage{unicode-math}
% unicode-math pulls amsmath in anyway; naming it makes the dependency explicit,
% since fragments use aligned/\dfrac and must typeset the same here as in the
% book, where gwbook.cls gets amsmath via physics.
\usepackage{amsmath}
\setmainfont{texgyreschola}[
  Extension=.otf, UprightFont=*-regular, ItalicFont=*-italic,
  BoldFont=*-bold, BoldItalicFont=*-bolditalic]
\setmathfont{texgyreschola-math.otf}

\usepackage{tikz}
% Same list, and for the same reason, as the preamble of gwbook.cls -- keep the
% two in step so a figure looks identical exported and on the page.
\usetikzlibrary{arrows.meta, calc, patterns.meta}

\begin{document}
\input{\FIG}
\end{document}
